% Extracto íntegro por residencia del propietario científico definitivo de masas.
\section{Raffinement explicite dans les tours \texorpdfstring{\(90/120\)}{90/120}}

L'existence du semi-groupe \(\langle90,120\rangle\) ne suffit pas à elle seule :
il faut construire l'application qui assigne des coefficients à une route enrichie. Soit
\[
 \mathcal H_{90,120}
 =\{90a+120b:a,b\in\mathbb N_0,\ (a,b)\ne(0,0)\}
 =\{90,120\}\cup\{30r:r\geq6\}
 =\{h_0<h_1<h_2<\cdots\}
\]
son énumération strictement croissante, où \(\mathbb N_0=\{0,1,2,\ldots\}\),
et soit \(s_{h_n}>0\) la normalisation déclarée du mode correspondant. Le
registre compatible d'une route
\(\Gamma\) fournit une suite de données entières finies
\[
 \mathcal L(\Gamma)=(\ell_0(\Gamma),\ell_1(\Gamma),\ldots),
\]
dans laquelle chaque \(\ell_n\) réunit le retour, la retenue, l'enroulement, la mémoire et la frontière
au raffinement \(n\). Pour fixer l'encodage sans ambiguïté, soit
\[
 \nu(z)=
 \begin{cases}2z,&z\geq0,\\-2z-1,&z<0,\end{cases}
 \qquad
 \pi_{\rm C}(u,v)=\frac{(u+v)(u+v+1)}2+v,
\]
et définissons \(\iota:\mathbb Z^d\to\mathbb N_0\) par application récursive de
\(\pi_{\rm C}\) à \(\nu(z_1),\ldots,\nu(z_d)\). Alors
\[
 q_n(\Gamma)
 =\frac{\iota(\ell_n(\Gamma))}
 {1+|\iota(\ell_n(\Gamma))|}\in(-1,1).
\]
La réalisation minimale de tour est
\[
 \eta_{h_n}(\Gamma)
 =\frac{3^{-(n+1)}q_n(\Gamma)}{s_{h_n}},
 \qquad
 \mathfrak T(\Gamma)
 =(\eta_h(\Gamma))_{h\in\mathcal H_{90,120}}.
\]

\begin{theorem}[Convergence, naturalité et cécité métrologique]
Pour toute route enrichie,
\[
 \sum_{h\in\mathcal H_{90,120}}|\eta_h(\Gamma)s_h|
 \leq\sum_{n\ge0}3^{-(n+1)}=\frac12.
\]
Si le raffinement prolonge le registre par préfixes, \(\mathfrak T\) est naturel
par rapport aux troncatures. De plus, aucune valeur de masse extérieure n'intervient
dans \(\mathfrak T\), et la queue après \(N\) modes satisfait
\[
 \sum_{n>N}|\eta_{h_n}s_{h_n}|
 \leq\frac{3^{-(N+1)}}2.
\]
\end{theorem}

\begin{proof}
On a \(|q_n|<1\) ; les deux bornes sont des sommes géométriques. La compatibilité
par préfixes conserve tous les coefficients déjà construits et ajoute ceux du
nouveau niveau, ce qui donne le carré de naturalité avec la troncature de
\(\mathcal E_{90,120}\). La formule ne contient que le registre, l'encodage
explicite et la normalisation modale déclarée.
\end{proof}

La construction précédente fixe une réalisation minimale explicite et convergente,
relative à l'encodage \(\iota\) et aux normalisations \(s_h\) déclarées.
Une dynamique qui modifie ses poids doit déclarer l'opérateur effectuant cette
modification et démontrer de nouveau la naturalité et la convergence ; elle ne peut pas
reconstruire les poids à partir d'un résidu expérimental.

\section{Opérateurs de fibre et loi bidirectionnelle}

Pour \(F_X=\mathscr S(X)\), soit
\[
 \mathcal H_X=\ell^2(F_X)\otimes V_{\rho_X}
\]
l'espace de réalisation, où \(V_{\rho_X}\) porte la représentation
interne correspondante. Chaque route \(\gamma\in F_X\) possède une signature
\[
 \sigma(\gamma)
 =(n_A,s_C,k_{\Delta_4},b_{120},b_\zeta)\in\mathbb Z^5
\]
et des coefficients de raffinement
\(\eta(\gamma)\in\mathcal E_{90,120}\). L'opérateur de base est
\[
 \widehat M_X^{(0)}
 =\sum_{\gamma\in F_X}
 m_{\rm clk}\,
 \chi_{\rm full}(\sigma(\gamma),\eta(\gamma))
 |\gamma\rangle\langle\gamma|\otimes I_{V_{\rho_X}}.
\]
L'échelle \(m_{\rm clk}\) descend de la branche absolue
action--horloge--masse, y compris la décennie \(10^{-34}\) ; elle n'est pas extraite d'une
masse expérimentale.

Les lecteurs \(K_D\) et \(K_\Omega\) induisent des opérateurs diagonaux exacts
\(B_X^D\) et \(B_X^\Omega\). Les deux orientations publient
\[
 \widehat M_X^{\beta,r}
 =R_{\rm act}^{B_X^r/2}\widehat M_X^{(0)}
  R_{\rm act}^{B_X^r/2},
 \qquad r\in\{D,\Omega\}.
\]
Les deux sont positives et commutent dans la base des routes. Par le théorème spectral
conjoint, il existe une unique mesure spectrale projective
\(E_X^{D,\Omega}\) sur \(\mathbb R_+^2\) telle que
\[
 \widehat M_X^{\beta,D}
 =\int x\,dE_X^{D,\Omega}(x,y),
 \qquad
 \widehat M_X^{\beta,\Omega}
 =\int y\,dE_X^{D,\Omega}(x,y).
\]
Dans la base des routes,
\[
 E_X^{D,\Omega}(A)|\gamma\rangle
 =\mathbf 1_A\!\left(m_D(\gamma),m_\Omega(\gamma)\right)
  |\gamma\rangle.
\]
C'est la sortie bidirectionnelle canonique : elle conserve la paire complète et ne
suppose pas d'avance une fonctionnelle scalaire.

Soit
\[
 \boldsymbol\beta_X(\gamma)
 =(m_D(\gamma),m_\Omega(\gamma)).
\]
La mesure image uniforme de la fibre est
\[
 \nu_X=(\boldsymbol\beta_X)_\#
 \left(\frac1{|F_X|}\sum_{\gamma\in F_X}\delta_\gamma\right)
 =\frac1{|F_X|}\sum_{\gamma\in F_X}
 \delta_{(m_D(\gamma),m_\Omega(\gamma))}.
\]
Sa transformée et ses moments mixtes sont
\[
 Z_X(s,t)=\int e^{sx+ty}\,d\nu_X(x,y),
 \qquad
 \mathfrak m_{p,q}(X)=\int x^py^q\,d\nu_X(x,y).
\]
Comme \(\nu_X\) est à support fini, \(Z_X\), ou de manière équivalente la famille
complète des moments mixtes, détermine le multiensemble conjoint des paires de
lecteurs à permutation des routes près. La mesure image est donc un
invariant complet de la paire opératorielle ; ce n'est pas une moyenne de masses.

Pour un canal \(P_{X,a}\), la mesure spectrale comprimée
\[
 F_{X,a}^{D,\Omega}(A)
 =\left.P_{X,a}E_X^{D,\Omega}(A)P_{X,a}\right|_{P_{X,a}\mathcal H_X}
\]
est une POVM sur le sous-espace préparé. Pour une densité \(\rho_X\)
portée par ce sous-espace, la distribution scalaire conjointe est
\[
 \mu_{X,a}^{D,\Omega}(A)
 =\operatorname{Tr}\!\left(
 \rho_X F_{X,a}^{D,\Omega}(A)
 \right).
\]
Toute fonctionnelle déclarée \(f:\mathbb R_+^2\to\mathbb R_+\) induit
\(\mu_{X,a}^f=f_\#\mu_{X,a}^{D,\Omega}\) sans modifier la PVM d'origine.

Si \(u:F_X\to F_Y\) est une bijection qui conserve le registre, l'orientation et
les deux lecteurs, l'unitaire \(U_u|\gamma\rangle=|u\gamma\rangle\) satisfait
\[
 U_uE_X^{D,\Omega}(A)=E_Y^{D,\Omega}(A)U_u,
 \qquad
 \nu_X=\nu_Y.
\]
La même égalité vaut pour les distributions physiques lorsque la densité et le
projecteur sont transportés par \(U_u\). Cette naturalité n'est affirmée que pour
les entrelaceurs effectivement construits ; une coïncidence de
cardinaux ne crée pas à elle seule une équivalence de fibres.

Lorsque la représentation physique exige une symétrie sous l'échange
\(D\leftrightarrow\Omega\), l'agrégat opératoriel minimal est la moyenne
géométrique
\[
 \widehat M_X^{\rm sym}
 =\widehat M_X^{\beta,D}\#\widehat M_X^{\beta,\Omega},
 \qquad
 A\#B
 =A^{1/2}\bigl(A^{-1/2}BA^{-1/2}\bigr)^{1/2}A^{1/2},
\]
sur le support positif commun. Dans le cas diagonal,
\[
 \widehat M_X^{\rm sym}
 =R_{\rm act}^{(B_X^D+B_X^\Omega)/4}
  \widehat M_X^{(0)}
  R_{\rm act}^{(B_X^D+B_X^\Omega)/4}.
\]
Cette formule est une conséquence de l'hypothèse explicite d'échange ; elle ne
remplace pas la mesure conjointe lorsque cette symétrie n'a pas été établie.

\begin{theorem}[Caractérisation de l'agrégat symétrique]
La moyenne géométrique est positive, symétrique, covariante par congruence pour
toute \(S\) inversible et autoduale :
\[
 A\#B=B\#A,
 \qquad
 S^*(A\#B)S=(S^*AS)\#(S^*BS),
 \qquad
 (A\#B)^{-1}=A^{-1}\#B^{-1}.
\]
De plus, \(A\#A=A\). Par conséquent, dans les secteurs où il existe un
entrelaceur physique qui échange les lecteurs, l'agrégat n'introduit pas de
paramètre et ne privilégie aucune orientation.
\end{theorem}

\begin{proof}
L'identité de congruence s'entend pour \(S\) inversible. Les identités
résultent du calcul fonctionnel de la racine positive de
\(A^{-1/2}BA^{-1/2}\). Dans le cas commutatif, elles se réduisent à
\(A\#B=(AB)^{1/2}\).
\end{proof}

Il existe en outre un agrégat scalaire canonique sous des axiomes distincts. Si
\(s_n\) est continu, invariant par permutation, multiplicatif coordonnée par
coordonnée et normalisé par \(s_n(cI)=c\), alors
\[
 s_n(T)=\det(T)^{1/n}.
\]
En effet, en passant aux logarithmes, la continuité et l'additivité rendent la
fonctionnelle linéaire sur les logarithmes des valeurs propres ; la symétrie égalise tous
les coefficients et la normalisation fixe leur valeur à \(1/n\). Pour la paire de
lecteurs, la symétrie d'échange donne l'invariant
\[
 S_{D,\Omega}
 =\sqrt{
 \operatorname{ndet}(T_D)\operatorname{ndet}(T_\Omega)},
 \qquad
 \operatorname{ndet}(T)=\det(T)^{1/n}.
\]
Ce résultat classifie la fibre complète sous les axiomes énoncés. Il ne s'
identifie pas automatiquement à une ligne de masse : cette identification exige
que le couplage physique adopte précisément cette fonctionnelle.

Pour la fibre électronique de rang un, la mesure conjointe est ponctuelle, les deux
lecteurs partagent la première composante \(80\), la compression est scalaire et l'on
retrouve la clôture bêta
\[
 m_e^\beta
 =0.5109989506900008086082571648\ldots\ {\rm MeV}.
\]
La valeur de base demeure une étape de construction ; elle ne remplace pas la clôture
bêta.

\section{Scalarisation par symétrie et couplage}

Soit \(G_X\) le groupe compact de symétrie non brisée de la fibre et
\(U_X:G_X\to\mathcal U(\mathcal H_X)\) sa représentation. L'espérance
conditionnelle de symétrie est définie par
\[
 \mathbb E_X(T)
 =\int_{G_X}U_X(g)TU_X(g)^*\,d\mu_X(g),
\]
où \(\mu_X\) est la mesure de Haar normalisée. L'opérateur physique de masse
avant l'ouverture des canaux est défini, lorsque le type d'état déclare une
fonctionnelle \(\Phi_X:\mathbb R_+^2\to\mathbb R_+\), par
\[
 \mathcal M_X
 =\mathbb E_X\!\left(
   \Phi_X(\widehat M_X^{\beta,D},\widehat M_X^{\beta,\Omega})
  \right).
\]
Le calcul fonctionnel s'effectue avec la mesure conjointe
\(E_X^{D,\Omega}\). Dans une représentation contenant un entrelaceur d'
échange des lecteurs, on peut prendre \(\Phi_X(x,y)=\sqrt{xy}\). Si le type
physique ne déclare pas \(\Phi_X\), la sortie complète reste la PVM conjointe ;
on n'invente pas de scalaire.

\begin{theorem}[Clôture scalaire, multiplet et distribution]
L'opérateur \(\mathcal M_X\) appartient au commutant
\(U_X(G_X)'\). Si
\[
 \mathcal H_X\simeq
 \bigoplus_\lambda V_\lambda\otimes\mathbb C^{m_\lambda},
\]
alors
\[
 \mathcal M_X=
 \bigoplus_\lambda I_{V_\lambda}\otimes M_{X,\lambda}.
\]
Dans un secteur irréductible de multiplicité \(1\), la restriction est un scalaire.
Dans un secteur réductible, la sortie exacte est le spectre des blocs
\(M_{X,\lambda}\), avec leurs multiplicités. Il n'existe pas d'indétermination :
il existe un codomaine distinct.
\end{theorem}

\begin{proof}
L'invariance de la mesure de Haar implique
\(U_X(h)\mathbb E_X(T)U_X(h)^*=\mathbb E_X(T)\) pour tout \(h\in G_X\).
La décomposition découle du théorème de complète réductibilité et du lemme de
Schur. Le cas irréductible produit un multiple de l'identité ; le cas
réductible conserve exactement les blocs de multiplicité.
\end{proof}

Soit \(P_{X,a}\) le projecteur déterminé par l'état préparé et le canal ou l'
appareil physique \(a\). La mesure spectrale observable est
\[
 \mu_{X,a}(B)=
 \langle\Psi_X,
 P_{X,a}E_{\mathcal M_X}(B)P_{X,a}\Psi_X\rangle.
\]
Il y a une ligne scalaire précisément lorsque le support spectral comprimé est
ponctuel :
\[
 E_{\mathcal M_X}(\{m_X\})P_{X,a}=P_{X,a},
 \qquad\text{de manière équivalente}\qquad
 (\mathcal M_X-m_X I)\,P_{X,a}=0.
\]
L'égalité isolée
\(P_{X,a}\mathcal M_X P_{X,a}=m_X P_{X,a}\) ne fixe qu'une espérance si le
sous-espace n'est pas réducteur et ne suffit pas à démontrer une ligne spectrale.
Si la compression conserve plusieurs valeurs propres, la sortie est un multiplet ou une
distribution. La remplacer par une moyenne numérique détruirait l'information
physique et généalogique.


